% Options for packages loaded elsewhere
\PassOptionsToPackage{unicode}{hyperref}
\PassOptionsToPackage{hyphens}{url}
%
\documentclass[
]{article}
\usepackage{amsmath,amssymb}
\usepackage{iftex}
\ifPDFTeX
  \usepackage[T1]{fontenc}
  \usepackage[utf8]{inputenc}
  \usepackage{textcomp} % provide euro and other symbols
\else % if luatex or xetex
  \usepackage{unicode-math} % this also loads fontspec
  \defaultfontfeatures{Scale=MatchLowercase}
  \defaultfontfeatures[\rmfamily]{Ligatures=TeX,Scale=1}
\fi
\usepackage{lmodern}
\ifPDFTeX\else
  % xetex/luatex font selection
\fi
% Explicit Unicode mapping for the pandoc-generated body under pdflatex
% (utf8 inputenc does not define Greek / math symbols by default).
\ifPDFTeX
\usepackage{newunicodechar}
\newunicodechar{Λ}{\ensuremath{\Lambda}}
\newunicodechar{Σ}{\ensuremath{\Sigma}}
\newunicodechar{Φ}{\ensuremath{\Phi}}
\newunicodechar{β}{\ensuremath{\beta}}
\newunicodechar{λ}{\ensuremath{\lambda}}
\newunicodechar{μ}{\ensuremath{\mu}}
\newunicodechar{ρ}{\ensuremath{\rho}}
\newunicodechar{τ}{\ensuremath{\tau}}
\newunicodechar{—}{---}
\newunicodechar{−}{\textminus}
\newunicodechar{≅}{\ensuremath{\cong}}
\newunicodechar{≈}{\ensuremath{\approx}}
\newunicodechar{⊗}{\ensuremath{\otimes}}
\newunicodechar{⊗}{\ensuremath{\otimes}}
\newunicodechar{×}{\ensuremath{\times}}
\newunicodechar{Γ}{\ensuremath{\Gamma}}
\newunicodechar{→}{\ensuremath{\to}}
% Table support (pandoc longtable output)
\usepackage{longtable,booktabs,array}
\usepackage{calc}
\fi
% Use upquote if available, for straight quotes in verbatim environments
\IfFileExists{upquote.sty}{\usepackage{upquote}}{}
\IfFileExists{microtype.sty}{% use microtype if available
  \usepackage[]{microtype}
  \UseMicrotypeSet[protrusion]{basicmath} % disable protrusion for tt fonts
}{}
\makeatletter
\@ifundefined{KOMAClassName}{% if non-KOMA class
  \IfFileExists{parskip.sty}{%
    \usepackage{parskip}
  }{% else
    \setlength{\parindent}{0pt}
    \setlength{\parskip}{6pt plus 2pt minus 1pt}}
}{% if KOMA class
  \KOMAoptions{parskip=half}}
\makeatother
\usepackage{xcolor}
\setlength{\emergencystretch}{3em} % prevent overfull lines
\providecommand{\tightlist}{%
  \setlength{\itemsep}{0pt}\setlength{\parskip}{0pt}}
\setcounter{secnumdepth}{-\maxdimen} % remove section numbering
\ifLuaTeX
  \usepackage{selnolig}  % disable illegal ligatures
\fi
\usepackage{bookmark}
\IfFileExists{xurl.sty}{\usepackage{xurl}}{} % add URL line breaks if available
\urlstyle{same}
\hypersetup{
  hidelinks,
  pdfcreator={LaTeX via pandoc}}

\author{}
\date{}

\begin{document}

\section{REM Specification v2.2}\label{rem-specification-v2.2}

\textbf{Version}: 2.2\\
\textbf{Date}: 2026-08-12\\
\textbf{Status}: Frozen (canonical candidate; D1/D2 re-verification, D3,
D4 open)\\
\textbf{Supersedes}: v2.1 (2026-08-12), v2.0 (2026-08-11)

\subsection{Summary of Changes from
v2.1}\label{summary-of-changes-from-v2.1}

v2.1 made the normalized instantaneous rate
\(\Gamma_F^{\mathrm{exact}}(0)\) an operational realization of the
dynamical cost and elevated a finite-time log-ratio functional to
primary status. \textbf{Phase D2.1 falsified the normalized rate as a
global variational objective}: on unrestricted TPS optimization the
denominator \(|Q_F\rho|^2\) can be driven to zero, sending
\(\Gamma_F \to \pm\infty\) and \(\Phi \to +\infty\) (a pure state can
always be approached by a product-state TPS). The v2.2 canonical
functional is therefore the \textbf{denominator-free (unnormalized)
instantaneous functional}

\[
\boxed{\,J_{\rm dyn}^{(0)}(F;\rho,\mathcal L) =
-\operatorname{Re}\langle Q_F\rho,\; Q_F\mathcal L(\rho)\rangle_{\mathrm{HS}}\,}
\]

with

\[
\boxed{\,\Phi(F;\lambda) = I_\rho(F) - \lambda\, J_{\rm dyn}^{(0)}(F)\,}
\]

as the canonical variational functional. The normalized rate is demoted
to a fixed-\((F)\) local diagnostic (Section 7). The finite-time
extension \(J_{\rm dyn}^{(\tau)}\) is repositioned as an
\textbf{operational finite-time extension / consistency diagnostic},
with the exact \(\tau\to 0\) limit
\(\lim_{\tau\to0}J_{\rm dyn}^{(\tau)}=J_{\rm dyn}^{(0)}\) as the main
consistency condition (Section 9). Phase D2.2-A/B validated this
structure: G1--G6 PASS (D2.2-A) and G7 \textbf{PASS with documented
finite-\(\tau\) crossover} (D2.2-B): for the Haar state,
\(\tau \ge 10^{-2}\) induces a genuine basin transition of \(F^*\) --- a
new result that defines the D3 phase.

Additional corrections relative to v2.1:

\begin{itemize}
\tightlist
\item
  v2.1's \(\Gamma_F^{\mathrm{exact}}\) carried a factor of 2 (from the
  squared norm convention). The C0 pre-registration defines the rate as
  \(-\mathrm{d}\log C_F/\mathrm{d}t\) with \(C_F\) the Frobenius
  \emph{norm}, giving \textbf{no factor of 2}; the numerical-derivative
  test and the master's D2.2 formula agree. Reverted in commit
  \texttt{fb08a3f}.
\item
  v2.1's finite-time functional was the log-ratio
  \(-(1/\tau)\log(C_F(\tau)/C_F(0))\). v2.2 freezes the difference ratio
  \(J_{\rm dyn}^{(\tau)} = -(C_F^2(\tau)-C_F^2(0))/(2\tau)\), which has
  the exact derivative-based \(\tau\to0\) limit above.
\end{itemize}

\begin{center}\rule{0.5\linewidth}{0.5pt}\end{center}

\subsection{1. Scope / Status}\label{scope-status}

REM treats the tensor-product structure of Hilbert space as a variable
structure selected by competing informational and dynamical constraints.
REM is a structural selection framework, not a modification of quantum
dynamics.

This specification freezes the \textbf{canonical variational functional}
of open-system REM after the D2.1 falsification:

\begin{itemize}
\tightlist
\item
  \textbf{Canonical}: \(J_{\rm dyn}^{(0)}(F;\rho,\mathcal L)\) (Section
  4) and \(\Phi(F;\lambda) = I_\rho - \lambda J_{\rm dyn}^{(0)}\)
  (Section 5).
\item
  \textbf{Prohibited as global objective}: the normalized rate
  \(\Gamma_F^{\mathrm{exact}}(0)\) (Section 7).
\item
  \textbf{Operational extension}: \(J_{\rm dyn}^{(\tau)}\) (Section 8)
  with the \(\tau\to0\) consistency relation (Section 9) and the
  admissible finite-\(\tau\) TPS crossover (Section 10).
\end{itemize}

Status: \textbf{frozen candidate}. Remaining validation gates: D1/D2
re-verification under the v2.2 canonical functional, D3 (timescale /
\(\tau\)-crossover), D4 (N-scaling) --- Section 12.

\subsection{2. Factorization domain (F)}\label{factorization-domain-f}

A factorization \(F\) is defined by a unitary \(U \in U(2^n)\) that
rotates the computational basis into a bipartite tensor-product
structure:

\[
\mathcal H \simeq \mathcal H_A \otimes \mathcal H_B,\qquad
n_A + n_B = n .
\]

The space of factorizations is the quotient

\[
\mathcal F = U(2^n)\big/ \mathrm{Stab}(A|B),
\]

where \(\mathrm{Stab}(A|B)\) is the stabilizer of the bipartition. The
quotient has real dimension \(D^2 - d_A^2 - d_B^2 + 1\) (the \(U(1)\)
kernel is included); for \(n=3\), \(n_A=2\): \(\dim \mathcal F = 45\).

The domain is \textbf{unrestricted}: \(F\) ranges over all TPS
rotations, not only contiguous site cuts. This is the setting in which
the D2.1 singularity arises (Section 7). Factorizations are compared
with the gauge-invariant distance
\(d_F(F_1,F_2) = \|P_{\mathcal A_{F_1}} - P_{\mathcal A_{F_2}}\|_F\)
(projector onto the local operator algebra), invariant under
\(V_A \otimes V_B\) gauge transformations.

\subsection{\texorpdfstring{3. Informational functional
(\(I_\rho(F)\))}{3. Informational functional (I\_\textbackslash rho(F))}}\label{informational-functional-i_rhof}

\[
I_\rho(F) = I(A:B)_\rho = S(\rho_A) + S(\rho_B) - S(\rho_{AB}),
\qquad
S(\rho) = -\mathrm{Tr}(\rho\log\rho).
\]

\(I_\rho\) is the mutual information of the (possibly mixed) state
\(\rho\) across the factorization \(F\). Units are nats or bits
depending on the logarithm base (the D-series numerics use \(\log_2\);
\(I \le \min(\log_2 d_A,
\log_2 d_B)\cdot 2\) for pure states). The frozen D-series protocol
evaluates \(I\) on the dominant eigenvector of the rotated state.

\subsection{\texorpdfstring{4. Canonical instantaneous dynamical
functional
(\(J_{\rm dyn}^{(0)}\))}{4. Canonical instantaneous dynamical functional (J\_\{\textbackslash rm dyn\}\^{}\{(0)\})}}\label{canonical-instantaneous-dynamical-functional-j_rm-dyn0}

\[
\boxed{\,
J_{\rm dyn}^{(0)}(F;\rho,\mathcal L)
= -\operatorname{Re}\langle Q_F\rho,\; Q_F\mathcal L(\rho)\rangle_{\mathrm{HS}}
\,}
\]

with

\begin{itemize}
\tightlist
\item
  \(Q_F = I - D_F\), where \(D_F\) is the dephasing projection onto the
  initial Schmidt basis of \(F\):
  \(D_F[\sigma] = \sum_k |b_k\rangle\langle b_k|
  \sigma |b_k\rangle\langle b_k|\),
\item
  \(\mathcal L\) the Lindblad Liouvillian of the fixed environment,
\item
  \(\langle A,B\rangle_{\mathrm{HS}} = \mathrm{Tr}(A^\dagger B)\).
\end{itemize}

\textbf{Name}: the \textbf{instantaneous dynamical rate functional
(signed)}. The functional is \textbf{signed} and must never be called a
``cost'' unconditionally: \(J_{\rm dyn}^{(0)} > 0\) when the initial
coherence \(C_F^2(t)=|Q_F\rho(t)|^2\) decays (penalty in \(\Phi\)),
\(J_{\rm dyn}^{(0)} < 0\) when it grows (reward in \(\Phi\)). By
Cauchy--Schwarz,

\[
|J_{\rm dyn}^{(0)}| \le |Q_F\rho|_2\, |Q_F\mathcal L(\rho)|_2 \le
|Q_F\rho|_2\, \|\mathcal L(\rho)\|_2,
\]

so \(J_{\rm dyn}^{(0)} \to 0\) smoothly as \(|Q_F\rho|^2 \to 0\):
\textbf{no denominator blow-up on the unrestricted TPS domain}. This is
the defining property that makes \(J_{\rm dyn}^{(0)}\) well-posed as a
global variational functional (D2.1 vs D2.2-A).

\subsection{\texorpdfstring{5. Canonical objective
(\(\Phi = I - \lambda J_{\rm dyn}^{(0)}\))}{5. Canonical objective (\textbackslash Phi = I - \textbackslash lambda J\_\{\textbackslash rm dyn\}\^{}\{(0)\})}}\label{canonical-objective-phi-i---lambda-j_rm-dyn0}

\[
\boxed{\,
\Phi(F;\lambda) = I_\rho(F) - \lambda\, J_{\rm dyn}^{(0)}(F;\rho,\mathcal L)
\,}
\]

Optimal factorization:

\[
F^*(\lambda) = \arg\max_{F\in\mathcal F} \Phi(F;\lambda).
\]

\textbf{Monotonic Tradeoff Theorem} (unchanged from v2.0, Section 8a):
for \(\lambda_2 > \lambda_1 \ge 0\),

\[
J_{\rm dyn}^{(0)}(F^*(\lambda_2)) \le J_{\rm dyn}^{(0)}(F^*(\lambda_1)),
\qquad
I_\rho(F^*(\lambda_2)) \le I_\rho(F^*(\lambda_1)),
\]

for a \textbf{fixed} environment \((\rho,\mathcal L)\). The theorem
holds for the signed functional with \(\lambda \ge 0\) (the proof
requires only non-negativity of \(\lambda\), not of the dynamical term).

\subsection{\texorpdfstring{6. Dimensional convention
(\([\lambda] = T\))}{6. Dimensional convention ({[}\textbackslash lambda{]} = T)}}\label{dimensional-convention-lambda-t}

\(\rho\) is dimensionless; \(\mathcal L\) has units \(T^{-1}\); hence

\[
[J_{\rm dyn}^{(0)}] = T^{-1},
\qquad
[\lambda] = T,
\qquad
[\lambda J_{\rm dyn}^{(0)}] = 1 .
\]

\(\lambda\) is the characteristic timescale that converts the
instantaneous rate into a dimensionless contribution to \(\Phi\). The
optional rescaled representation \(\tilde J = \tau_0 J\),
\(\tilde\lambda = \lambda/\tau_0\) (with
\(\Phi = I - \tilde\lambda \tilde J\)) is mathematically equivalent and
introduces no new physics; the canonical form keeps \([\lambda]=T\) with
no reference time \(\tau_0\).

\subsection{\texorpdfstring{7. Normalized \(\Gamma_F\):
local-diagnostic-only
status}{7. Normalized \textbackslash Gamma\_F: local-diagnostic-only status}}\label{normalized-gamma_f-local-diagnostic-only-status}

The normalized instantaneous rate

\[
\Gamma_F^{\mathrm{exact}}(0)
= -\frac{\operatorname{Re}\langle Q_F\rho,\; Q_F\mathcal L(\rho)\rangle_{\mathrm{HS}}}
{|Q_F\rho|_2^2}
\]

is \textbf{retained but demoted}:

\begin{quote}
\textbf{fixed-\((F)\) local diagnostic only; prohibited as an
unrestricted global TPS variational objective.}
\end{quote}

Rationale (Phase D2.1, commit \texttt{0ad8840}): on unrestricted TPS
optimization the optimizer drives \(|Q_F\rho|^2 \to 0\), which sends
\(\Gamma_F \to
\pm\infty\) and \(\Phi \to +\infty\) along a non-physical direction.
Observed evidence:

\begin{itemize}
\tightlist
\item
  \(\operatorname{corr}(\Phi, \log_{10} C_F(0)^2) = -0.999\)
  (near-perfect anti-correlation),
\item
  \(C_F(0)^2 \to 6\times10^{-5}\) drives \(\Gamma \to -222\),
  \(\Phi \to 44\),
\item
  best solution Schmidt spectrum \(p = [0.9997, 0.0003]\) (near-product
  collapse).
\end{itemize}

Because every pure state can be approached by a product-state TPS, the
denominator can vanish; the normalized rate is therefore well-posed only
for a \textbf{fixed} \(F\) (where \(|Q_F\rho|^2\) is bounded away from
zero) as a local diagnostic of structural decoherence. The v2.1
elevation of \(\Gamma_F\) to an operational canonical cost is
\textbf{abandoned} on this falsification.

\subsection{\texorpdfstring{8. Finite-time extension
(\(J_{\rm dyn}^{(\tau)}\))}{8. Finite-time extension (J\_\{\textbackslash rm dyn\}\^{}\{(\textbackslash tau)\})}}\label{finite-time-extension-j_rm-dyntau}

\[
\boxed{\,
J_{\rm dyn}^{(\tau)}(F;\rho,\mathcal L)
= -\frac{C_F^2(\tau) - C_F^2(0)}{2\tau}
\,}
\]

with \(C_F^2(t) = |Q_F \rho(t)|_2^2\),
\(\rho(t) = e^{t\mathcal L}[\rho]\), and \(Q_F\) built from the
\textbf{fixed} \(t=0\) Schmidt basis of \(F\).

\textbf{Status}: \textbf{operational finite-time extension / consistency
diagnostic} --- separate from the canonical functional, not an
alternative to it. It measures the average rate of squared-coherence
change over \([0,\tau]\) and is the natural finite-window probe of
structural stability.

\subsection{\texorpdfstring{9. \(\tau\to0\) consistency
relation}{9. \textbackslash tau\textbackslash to0 consistency relation}}\label{tauto0-consistency-relation}

For fixed \(F\) (hence fixed \(Q_F\)),

\[
\frac{\mathrm{d}}{\mathrm{d}t} C_F^2(t)\bigg|_{t=0}
= 2\operatorname{Re}\langle Q_F\rho,\; Q_F\mathcal L(\rho)\rangle_{\mathrm{HS}},
\]

so

\[
\boxed{\,
\lim_{\tau\to0} J_{\rm dyn}^{(\tau)}(F;\rho,\mathcal L)
= J_{\rm dyn}^{(0)}(F;\rho,\mathcal L)
\,}
\]

is the \textbf{main consistency condition of v2.2}: the finite-time
extension reproduces the canonical instantaneous functional in the
\(\tau\to0\) limit. Verified numerically (D2.2-B, G7-1): deviation
scales linearly in \(\tau\) (\(|\Delta J| \approx |J_0|\,\tau\)), with
\(|\Delta J| \le 2.5\times10^{-3}\) at \(\tau = 10^{-3}\) for all four
state classes.

\subsection{\texorpdfstring{10. Finite-\(\tau\) TPS crossover
(admissible)}{10. Finite-\textbackslash tau TPS crossover (admissible)}}\label{finite-tau-tps-crossover-admissible}

The consistency relation does \textbf{not} imply that \(F^*\) is
\(\tau\)-independent at finite \(\tau\):

\begin{quote}
\textbf{Finite \(\tau\) need not preserve the instantaneous optimum;
\(\tau\) may induce genuine transitions between competing TPS basins.}
\end{quote}

Evidence (D2.2-B, Haar state): for \(\tau \le 3\times10^{-3}\) all seeds
stay in the instantaneous basin (\(d_F < 0.06\)); for
\(\tau \ge 10^{-2}\) all seeds move to a different basin (\(I\): 1.97 →
2.00, \(J_{\rm dyn}^{(\tau)}\): 0.60 → 0.44, \(\Phi\) up to +4.3\% at
\(\tau = 0.1\)). The shift is robust (6/6 seeds), monotone in \(\tau\),
and free of pathology (no singularity, no product collapse, no seed
instability). It is read as a \textbf{genuine structural crossover
induced by finite-time coarse-graining}, not as noise and not as a
failure of the canonical functional. This is the defining phenomenon of
the D3 phase (Section 12).

\subsection{11. Known numerical evidence (D-series,
2026-08-12)}\label{known-numerical-evidence-d-series-2026-08-12}

All runs use the frozen protocol: asymmetric-XY chain (\(J_{12}=1.5\),
\(J_{23}=0.6\), \(h=0.2\)), pure dephasing \(\gamma=(0.5,1.0,2.0)\),
\(\lambda=0.2\), Adam 200 steps, \(\mathrm{lr}=0.01\), seeds
\(20260813\ldots\), quotient \(\dim\mathcal F=45\), \(n_A=2\). Suite: 74
passed + 1 xfailed.

{\def\LTcaptype{none} % do not increment counter
\begin{longtable}[]{@{}
  >{\raggedright\arraybackslash}p{(\linewidth - 4\tabcolsep) * \real{0.3333}}
  >{\raggedright\arraybackslash}p{(\linewidth - 4\tabcolsep) * \real{0.3333}}
  >{\raggedright\arraybackslash}p{(\linewidth - 4\tabcolsep) * \real{0.3333}}@{}}
\toprule\noalign{}
\begin{minipage}[b]{\linewidth}\raggedright
Phase
\end{minipage} & \begin{minipage}[b]{\linewidth}\raggedright
Result
\end{minipage} & \begin{minipage}[b]{\linewidth}\raggedright
Verdict
\end{minipage} \\
\midrule\noalign{}
\endhead
\bottomrule\noalign{}
\endlastfoot
D1 (Hamiltonian generality, Γ-based) & 5/5 families show non-trivial
\(F^*\) & preliminary (to be re-run under \(J_{\rm dyn}^{(0)}\)) \\
D2 (state generality, Γ-based) & ground/mixed/thermal PASS-candidate;
Haar anomaly (\(\Phi=49.1\)) & re-scoped to D2.1 \\
\textbf{D2.1} (Haar singularity audit, \texttt{0ad8840}) & normalization
singularity confirmed & \textbf{falsification} → Section 7 \\
\textbf{D2.2-A} (unnormalized instantaneous, \texttt{24db549}) & G1--G6
\textbf{PASS} & \(J_{\rm dyn}^{(0)}\) well-posed \\
\textbf{D2.2-B} (finite-time, \texttt{2fce54d}) & G7 \textbf{PASS with
documented finite-\(\tau\) crossover} & extension + Section 10 \\
\end{longtable}
}

D2.2-A numbers (4 states × 6 seeds): best \(\Phi\) = ground 1.8995, Haar
1.8594, mixed 1.9845, thermal 1.9465; Haar 30-seed best 1.8595 / std
0.0077 (vs 25.71 / 5.96 under the normalized objective); min
\(C_F^2 = 0.096\) (never approaches zero); Schmidt
\(p_{\max} = 0.51\)--\(0.58\) (no collapse); small- perturbation
continuity max \(|\Delta\Phi| = 7.7\times10^{-5}\).

D2.2-B numbers (τ sweep
\(\{10^{-3}, 3\times10^{-3}, 10^{-2}, 3\times10^{-2},
10^{-1}\}\)): G7-1 exact \(\tau\to0\) limit (above); G7-2 same-basin for
all states at \(\tau\le3\times10^{-3}\) and for ground/mixed/thermal at
all \(\tau\) (cross-eval gaps \(\le 0.005\)), Haar crossover at
\(\tau\ge10^{-2}\) (Section 10); G7-3 state ranking (mixed
\textgreater{} thermal \textgreater{} ground \textgreater{} Haar)
preserved at all \(\tau\); G7-4 min \(C_F^2(0) \ge 0.09\),
\(p_{\max} \le 0.58\), std \(\le 0.009\) at all
\((\mathrm{state},\tau)\); Haar 30-seed at \(\tau=0.1\): best 1.9135 /
std 0.0011.

\subsection{12. Open validation gates}\label{open-validation-gates}

\begin{enumerate}
\def\labelenumi{\arabic{enumi}.}
\tightlist
\item
  \textbf{D1 re-verification}: Hamiltonian generality (5 families) under
  the v2.2 canonical functional \(J_{\rm dyn}^{(0)}\).
\item
  \textbf{D2 re-verification}: state generality (ground / Haar / mixed /
  thermal) under \(J_{\rm dyn}^{(0)}\) (D2.2-A is the Haar-focused
  audit; the full 4-state protocol is the re-run).
\item
  \textbf{D3 Timescale}: promoted from ``how to determine \(\tau\)'' to
  \textbf{\(\tau\)-crossover analysis}: locate the basin-transition
  scale \(\tau_c\) (for Haar the switch lies between \(3\times10^{-3}\)
  and \(10^{-2}\)), characterize its sharpness, and connect it to
  environment timescales (e.g.~Liouvillian gap).
\item
  \textbf{D4 N-scaling}: \(N=4,5,\ldots\) behavior of \(F^*\),
  \(J_{\rm dyn}^{(0)}\), and the crossover.
\item
  Existing roadmap gates 3a/3b (λ/τ determination), Gate 4 (predictive
  hold-out), Gate 5 (generality) remain OPEN.
\end{enumerate}

\begin{center}\rule{0.5\linewidth}{0.5pt}\end{center}

\subsection{Version History}\label{version-history}

\begin{itemize}
\tightlist
\item
  \textbf{v1.0} (2026-04-25): initial formulation with closed-system
  surrogate \(C_H^{\mathrm{closed}} = \langle H_{\partial F}^2\rangle\).
\item
  \textbf{v1.1} (2026-07-18): monotonic tradeoff theorem;
  sign-convention fix.
\item
  \textbf{v2.0} (2026-08-11): open-system canonical form
  \(\Phi = I - \lambda C_{\mathrm{dyn}}\) with \([\lambda]=T\); first
  realization \(C_\Gamma^{(0)} = \Gamma_F^{\mathrm{exact}}(0)\);
  \(C_H^{\mathrm{closed}}\) downgraded to surrogate; signed dynamical
  functional named.
\item
  \textbf{v2.1} (2026-08-12): timescale-dependent formulation;
  finite-time log-ratio functional made primary; \(\Gamma_F\) retained
  as canonical cost. \textbf{Superseded by v2.2} on the D2.1
  falsification (Section 7): the normalized rate is singular on
  unrestricted TPS and is demoted to a local diagnostic; the
  unnormalized instantaneous functional is frozen as canonical (Sections
  4--5); the finite-time difference ratio is frozen as the operational
  extension (Section 8).
\item
  \textbf{v2.2} (2026-08-12): this document. Canonical
  \(J_{\rm dyn}^{(0)} = -\operatorname{Re}\langle Q_F\rho, Q_F\mathcal L(\rho)\rangle\);
  \(\Phi = I - \lambda J_{\rm dyn}^{(0)}\); \([J] = T^{-1}\),
  \([\lambda]=T\); G7 recorded as PASS with documented finite-\(\tau\)
  crossover.
\end{itemize}

\subsection{References}\label{references}

\begin{itemize}
\tightlist
\item
  D-series scripts: \texttt{analysis/d0\_protocol.py},
  \texttt{analysis/d1\_generality.py},
  \texttt{analysis/d2\_state\_generality.py},
  \texttt{analysis/d2\_1\_haar\_singularity\_audit.py},
  \texttt{analysis/d2\_2a\_unnormalized\_instantaneous.py},
  \texttt{analysis/d2\_2b\_finite\_time.py},
  \texttt{analysis/d2\_2b\_gates.py}.
\item
  Numerical results: \texttt{analysis\_output/d2\_*.json},
  \texttt{analysis\_output/d2\_2b\_*.json}.
\item
  Roadmap: \texttt{REM\_next\_steps.md} (D0--D2.2-B status).
\item
  Zenodo: DOI 10.5281/zenodo.21880505 (REM\_lambda v5); v6 correction
  plan will include D2.1--D2.2 and Spec v2.2.
\end{itemize}

\end{document}
